\documentclass[UTF8,a4paper,14pt]{ctexart}

\usepackage{geometry}
\usepackage{amsmath,amssymb}
\usepackage{longtable,booktabs,array}
\usepackage{xcolor}
\usepackage{listings}
\usepackage{hyperref}
\usepackage{enumitem}
\usepackage{fancyhdr}
\hypersetup{
    colorlinks=false, % 不显示链接颜色
    pdfborder={0 0 0}  % 设置边框宽度为 0
}
\geometry{left=2.2cm,right=2.2cm,top=2.2cm,bottom=2.2cm}
\setlist[itemize]{leftmargin=2em}
\setlist[enumerate]{leftmargin=2em}
\pagestyle{fancy}
\fancyhf{}
\lhead{百度之星算法速查}
\rhead{Python 模板}
\cfoot{\thepage}

\definecolor{codebg}{RGB}{248,248,248}
\definecolor{keyword}{RGB}{0,92,175}
\definecolor{comment}{RGB}{80,130,80}
\definecolor{string}{RGB}{160,80,40}

\lstset{
  language=Python,
  backgroundcolor=\color{codebg},
  basicstyle=\ttfamily\small,
  keywordstyle=\color{keyword}\bfseries,
  commentstyle=\color{comment},
  stringstyle=\color{string},
  showstringspaces=false,
  breaklines=true,
  frame=single,
  tabsize=4,
  columns=fullflexible,
  keepspaces=true
}

\title{\textbf{百度之星 / 算法竞赛常用算法总结}\\
\large 算法流程、Python 解题模板与题型识别}
\author{整理版}
\date{}

\begin{document}
\maketitle
\tableofcontents
\newpage

\section{总览：如何识别题型}

\subsection{常见题型与算法映射}

\begin{longtable}{p{0.25\textwidth}p{0.33\textwidth}p{0.32\textwidth}}
\toprule
\textbf{题面信号} & \textbf{优先考虑} & \textbf{典型关键词} \\
\midrule
求区间和、区间次数、区间最值 & 前缀和、差分、树状数组、线段树、ST 表 & 多次询问、区间修改、静态 RMQ \\
最小化最大值、最大化最小值 & 二分答案 & 最少时间、最大距离、最小代价、能否完成 \\
连续子数组、双端移动 & 双指针、滑动窗口、前缀和 + 哈希 & 最长、最短、恰好/至少/至多 \\
维护动态集合最值 & 堆、平衡思想、懒删除 & 第 $k$ 大、实时中位数、优先级 \\
连通性、合并集合 & 并查集、DFS/BFS & 联通块、朋友关系、合并、判断同组 \\
最短路、路径代价 & Dijkstra、BFS、Floyd & 路径、代价、距离、边权 \\
有向依赖关系 & 拓扑排序、DAG DP & 先后顺序、课程、任务依赖 \\
选择若干物品最优 & 背包 DP、贪心、排序 & 容量、价值、收益、限制 \\
字符串匹配 & KMP、哈希、Trie、Z 函数 & 子串、前缀、模式串、字典树 \\
取模计数 & 组合数学、快速幂、DP & $10^9+7$、方案数、排列组合 \\
博弈轮流操作 & SG 函数、奇偶性、Nim & 两人轮流、必胜、必败 \\
坐标很大但点少 & 离散化、扫描线 & 坐标到 $10^9$、事件、覆盖 \\
\bottomrule
\end{longtable}

\subsection{Python 竞赛基本框架}

\begin{lstlisting}
import sys
from collections import deque, defaultdict, Counter
from heapq import heappush, heappop
from bisect import bisect_left, bisect_right
from math import gcd, isqrt

input = sys.stdin.readline
MOD = 10 ** 9 + 7
INF = 10 ** 30

def solve():
    n = int(input())
    a = list(map(int, input().split()))
    # write solution here

if __name__ == "__main__":
    solve()
\end{lstlisting}

\subsection{复杂度速查}

\begin{itemize}
  \item $n \le 20$：可考虑状态压缩、枚举子集、Meet in the Middle。
  \item $n \le 10^3$：可考虑 $O(n^2)$ DP 或图算法。
  \item $n \le 10^5$：一般需要 $O(n\log n)$ 或 $O(n)$。
  \item $n \le 10^6$：一般需要 $O(n)$，并注意输入输出效率。
  \item 多组测试时，总复杂度看 $\sum n$，不要只看单组 $n$。
\end{itemize}

\section{基础算法}

\subsection{排序与自定义排序}

\textbf{流程：}
\begin{enumerate}
  \item 找出比较关键字。
  \item 判断升序还是降序。
  \item 若有多个条件，使用元组排序。
\end{enumerate}

\textbf{适用题型：}
区间排序、贪心选择、按结束时间排序、按收益排序、事件扫描。

\begin{lstlisting}
a.sort()                       # ascending
a.sort(reverse=True)           # descending
items.sort(key=lambda x: (x[0], -x[1]))
\end{lstlisting}

\subsection{前缀和}

\textbf{核心思想：}
用 $pre[i]$ 表示前 $i$ 个元素之和，则区间 $[l,r]$ 的和为
\[
sum(l,r)=pre[r+1]-pre[l].
\]

\textbf{算法流程：}
\begin{enumerate}
  \item 建立长度为 $n+1$ 的前缀数组。
  \item $pre[i+1]=pre[i]+a[i]$。
  \item 查询时用差值快速得到区间信息。
\end{enumerate}

\textbf{适用题型：}
多次区间求和、区间统计、子数组和、二维矩阵区域和。

\begin{lstlisting}
def prefix_sum(a):
    pre = [0]
    for x in a:
        pre.append(pre[-1] + x)
    return pre

def range_sum(pre, l, r):
    return pre[r + 1] - pre[l]
\end{lstlisting}

\subsection{二维前缀和}

二维前缀和适合矩阵子矩形求和。

\[
S[x_2][y_2]-S[x_1][y_2]-S[x_2][y_1]+S[x_1][y_1]
\]

\begin{lstlisting}
def build_2d_prefix(mat):
    n, m = len(mat), len(mat[0])
    s = [[0] * (m + 1) for _ in range(n + 1)]
    for i in range(n):
        row_sum = 0
        for j in range(m):
            row_sum += mat[i][j]
            s[i + 1][j + 1] = s[i][j + 1] + row_sum
    return s

def rect_sum(s, x1, y1, x2, y2):
    # rectangle [x1, x2) * [y1, y2)
    return s[x2][y2] - s[x1][y2] - s[x2][y1] + s[x1][y1]
\end{lstlisting}

\subsection{差分数组}

\textbf{核心思想：}
若要对区间 $[l,r]$ 加 $v$，只需：
\[
diff[l]+=v,\quad diff[r+1]-=v.
\]

最后对差分数组做一次前缀和即可还原。

\textbf{适用题型：}
多次区间加、最终数组、区间覆盖次数、离线修改。

\begin{lstlisting}
def apply_range_add(n, ops):
    diff = [0] * (n + 1)
    for l, r, v in ops:
        diff[l] += v
        if r + 1 < n:
            diff[r + 1] -= v
    ans = [0] * n
    cur = 0
    for i in range(n):
        cur += diff[i]
        ans[i] = cur
    return ans
\end{lstlisting}

\subsection{离散化}

\textbf{适用条件：}
坐标值很大，但实际出现的点较少。

\textbf{流程：}
\begin{enumerate}
  \item 收集所有有用坐标。
  \item 排序去重。
  \item 用字典或二分把原坐标映射到较小下标。
\end{enumerate}

\begin{lstlisting}
def compress(vals):
    xs = sorted(set(vals))
    mp = {x: i for i, x in enumerate(xs)}
    return xs, mp
\end{lstlisting}

\subsection{二分查找}

\textbf{左边界模板：}
寻找第一个满足 $a[i]\ge x$ 的位置。

\begin{lstlisting}
def lower_bound(a, x):
    l, r = 0, len(a)
    while l < r:
        mid = (l + r) // 2
        if a[mid] >= x:
            r = mid
        else:
            l = mid + 1
    return l
\end{lstlisting}

\subsection{二分答案}

\textbf{核心判断：}
答案具有单调性。即某个值可行，则更大或更小的一侧也可行。

\textbf{流程：}
\begin{enumerate}
  \item 确定答案范围 $[L,R]$。
  \item 写 \texttt{check(mid)} 判断 $mid$ 是否可行。
  \item 根据可行性收缩区间。
\end{enumerate}

\textbf{适用题型：}
最小最大值、最大最小值、最短时间、最少代价、容量问题。

\begin{lstlisting}
def binary_answer_left(lo, hi, check):
    # find minimum x such that check(x) is True
    while lo < hi:
        mid = (lo + hi) // 2
        if check(mid):
            hi = mid
        else:
            lo = mid + 1
    return lo

def binary_answer_right(lo, hi, check):
    # find maximum x such that check(x) is True
    while lo < hi:
        mid = (lo + hi + 1) // 2
        if check(mid):
            lo = mid
        else:
            hi = mid - 1
    return lo
\end{lstlisting}

\section{双指针、滑动窗口与单调结构}

\subsection{双指针}

\textbf{适用题型：}
有序数组两数和、区间满足某条件、去重、合并。

\begin{lstlisting}
def two_sum_sorted(a, target):
    l, r = 0, len(a) - 1
    while l < r:
        s = a[l] + a[r]
        if s == target:
            return l, r
        elif s < target:
            l += 1
        else:
            r -= 1
    return -1, -1
\end{lstlisting}

\subsection{滑动窗口}

\textbf{适用条件：}
窗口扩大或缩小时，性质能被维护，常见于正数数组、字符计数。

\begin{lstlisting}
def min_len_at_least_s(a, S):
    ans = 10 ** 18
    cur = 0
    l = 0
    for r, x in enumerate(a):
        cur += x
        while cur >= S:
            ans = min(ans, r - l + 1)
            cur -= a[l]
            l += 1
    return ans if ans < 10 ** 18 else -1
\end{lstlisting}

\subsection{单调栈}

\textbf{适用题型：}
下一个更大/更小元素、柱状图最大矩形、贡献法计算。

\begin{lstlisting}
def next_greater(a):
    n = len(a)
    ans = [-1] * n
    st = []
    for i, x in enumerate(a):
        while st and a[st[-1]] < x:
            ans[st.pop()] = i
        st.append(i)
    return ans
\end{lstlisting}

\subsection{单调队列}

\textbf{适用题型：}
固定长度滑动窗口最大值/最小值，DP 优化。

\begin{lstlisting}
from collections import deque

def sliding_window_max(a, k):
    q = deque()
    ans = []
    for i, x in enumerate(a):
        while q and q[0] <= i - k:
            q.popleft()
        while q and a[q[-1]] <= x:
            q.pop()
        q.append(i)
        if i >= k - 1:
            ans.append(a[q[0]])
    return ans
\end{lstlisting}

\section{贪心}

\subsection{贪心的判断方法}

\textbf{常见可贪心信号：}
\begin{itemize}
  \item 每一步局部选择不会影响后续最优性。
  \item 可以排序后选择。
  \item 可以用交换论证证明。
  \item 题目要求最大数量、最早结束、最小代价。
\end{itemize}

\subsection{区间调度}

选择最多不重叠区间，按右端点升序排序。

\begin{lstlisting}
def max_non_overlap(intervals):
    intervals.sort(key=lambda x: x[1])
    ans = 0
    last_end = -10 ** 30
    for l, r in intervals:
        if l >= last_end:
            ans += 1
            last_end = r
    return ans
\end{lstlisting}

\subsection{区间合并}

\begin{lstlisting}
def merge_intervals(intervals):
    intervals.sort()
    res = []
    for l, r in intervals:
        if not res or l > res[-1][1]:
            res.append([l, r])
        else:
            res[-1][1] = max(res[-1][1], r)
    return res
\end{lstlisting}

\section{数据结构}

\subsection{堆}

\textbf{适用题型：}
动态最小值、动态最大值、Top-K、合并有序序列、Dijkstra。

\begin{lstlisting}
from heapq import heappush, heappop

pq = []
heappush(pq, (cost, node))
cost, node = heappop(pq)

# max heap
heappush(pq, -x)
x = -heappop(pq)
\end{lstlisting}

\subsection{懒删除堆}

\begin{lstlisting}
from heapq import heappush, heappop
from collections import Counter

class LazyHeap:
    def __init__(self):
        self.h = []
        self.del_cnt = Counter()

    def push(self, x):
        heappush(self.h, x)

    def erase(self, x):
        self.del_cnt[x] += 1

    def _clean(self):
        while self.h and self.del_cnt[self.h[0]]:
            x = heappop(self.h)
            self.del_cnt[x] -= 1

    def top(self):
        self._clean()
        return self.h[0] if self.h else None

    def pop(self):
        self._clean()
        return heappop(self.h)
\end{lstlisting}

\subsection{并查集}

\textbf{适用题型：}
连通性、合并集合、最小生成树、动态加边。

\begin{lstlisting}
class DSU:
    def __init__(self, n):
        self.fa = list(range(n))
        self.sz = [1] * n

    def find(self, x):
        while self.fa[x] != x:
            self.fa[x] = self.fa[self.fa[x]]
            x = self.fa[x]
        return x

    def union(self, a, b):
        ra, rb = self.find(a), self.find(b)
        if ra == rb:
            return False
        if self.sz[ra] < self.sz[rb]:
            ra, rb = rb, ra
        self.fa[rb] = ra
        self.sz[ra] += self.sz[rb]
        return True
\end{lstlisting}

\subsection{树状数组}

\textbf{功能：}
单点修改、前缀查询。也可配合差分做区间加。

\begin{lstlisting}
class BIT:
    def __init__(self, n):
        self.n = n
        self.t = [0] * (n + 1)

    def add(self, i, v):
        i += 1
        while i <= self.n:
            self.t[i] += v
            i += i & -i

    def sum(self, i):
        # sum of [0, i]
        i += 1
        res = 0
        while i > 0:
            res += self.t[i]
            i -= i & -i
        return res

    def range_sum(self, l, r):
        if r < l:
            return 0
        return self.sum(r) - (self.sum(l - 1) if l else 0)
\end{lstlisting}

\subsection{线段树}

\textbf{适用题型：}
区间查询、区间修改、动态维护区间和/最值/最大子段和等。

\begin{lstlisting}
class SegTree:
    def __init__(self, a):
        self.n = len(a)
        self.sum = [0] * (4 * self.n)
        self.lazy = [0] * (4 * self.n)
        self._build(1, 0, self.n - 1, a)

    def _build(self, o, l, r, a):
        if l == r:
            self.sum[o] = a[l]
            return
        m = (l + r) // 2
        self._build(o * 2, l, m, a)
        self._build(o * 2 + 1, m + 1, r, a)
        self.sum[o] = self.sum[o * 2] + self.sum[o * 2 + 1]

    def _apply(self, o, l, r, v):
        self.sum[o] += (r - l + 1) * v
        self.lazy[o] += v

    def _push(self, o, l, r):
        if self.lazy[o]:
            m = (l + r) // 2
            v = self.lazy[o]
            self._apply(o * 2, l, m, v)
            self._apply(o * 2 + 1, m + 1, r, v)
            self.lazy[o] = 0

    def add(self, L, R, v):
        self._add(1, 0, self.n - 1, L, R, v)

    def _add(self, o, l, r, L, R, v):
        if L <= l and r <= R:
            self._apply(o, l, r, v)
            return
        self._push(o, l, r)
        m = (l + r) // 2
        if L <= m:
            self._add(o * 2, l, m, L, R, v)
        if R > m:
            self._add(o * 2 + 1, m + 1, r, L, R, v)
        self.sum[o] = self.sum[o * 2] + self.sum[o * 2 + 1]

    def query(self, L, R):
        return self._query(1, 0, self.n - 1, L, R)

    def _query(self, o, l, r, L, R):
        if L <= l and r <= R:
            return self.sum[o]
        self._push(o, l, r)
        m = (l + r) // 2
        res = 0
        if L <= m:
            res += self._query(o * 2, l, m, L, R)
        if R > m:
            res += self._query(o * 2 + 1, m + 1, r, L, R)
        return res
\end{lstlisting}

\subsection{ST 表}

\textbf{适用题型：}
静态区间最值查询，不能修改。

\begin{lstlisting}
from math import log2

class SparseTable:
    def __init__(self, a, func=min):
        self.func = func
        n = len(a)
        K = n.bit_length()
        self.st = [a[:]]
        j = 1
        while (1 << j) <= n:
            prev = self.st[-1]
            length = 1 << (j - 1)
            cur = [
                func(prev[i], prev[i + length])
                for i in range(n - (1 << j) + 1)
            ]
            self.st.append(cur)
            j += 1

    def query(self, l, r):
        # inclusive [l, r]
        k = (r - l + 1).bit_length() - 1
        return self.func(self.st[k][l], self.st[k][r - (1 << k) + 1])
\end{lstlisting}

\section{图论}

\subsection{建图模板}

\begin{lstlisting}
def build_graph(n, edges, directed=False):
    g = [[] for _ in range(n)]
    for u, v in edges:
        g[u].append(v)
        if not directed:
            g[v].append(u)
    return g
\end{lstlisting}

\subsection{BFS}

\textbf{适用题型：}
无权最短路、网格最短路、层次遍历。

\begin{lstlisting}
from collections import deque

def bfs(g, s):
    n = len(g)
    dist = [-1] * n
    dist[s] = 0
    q = deque([s])
    while q:
        u = q.popleft()
        for v in g[u]:
            if dist[v] == -1:
                dist[v] = dist[u] + 1
                q.append(v)
    return dist
\end{lstlisting}

\subsection{DFS}

\textbf{适用题型：}
连通块、树 DP、回溯、拓扑辅助。

\begin{lstlisting}
def dfs_iter(g, s):
    n = len(g)
    vis = [False] * n
    st = [s]
    vis[s] = True
    order = []
    while st:
        u = st.pop()
        order.append(u)
        for v in g[u]:
            if not vis[v]:
                vis[v] = True
                st.append(v)
    return order
\end{lstlisting}

\subsection{拓扑排序}

\textbf{适用题型：}
有向无环图、课程安排、任务依赖、DAG DP。

\begin{lstlisting}
from collections import deque

def topo_sort(n, edges):
    g = [[] for _ in range(n)]
    indeg = [0] * n
    for u, v in edges:
        g[u].append(v)
        indeg[v] += 1

    q = deque(i for i in range(n) if indeg[i] == 0)
    order = []
    while q:
        u = q.popleft()
        order.append(u)
        for v in g[u]:
            indeg[v] -= 1
            if indeg[v] == 0:
                q.append(v)

    return order if len(order) == n else []
\end{lstlisting}

\subsection{Dijkstra}

\textbf{适用条件：}
边权非负。

\begin{lstlisting}
from heapq import heappush, heappop

def dijkstra(n, g, s):
    # g[u] = [(v, w), ...]
    INF = 10 ** 30
    dist = [INF] * n
    dist[s] = 0
    pq = [(0, s)]
    while pq:
        d, u = heappop(pq)
        if d != dist[u]:
            continue
        for v, w in g[u]:
            nd = d + w
            if nd < dist[v]:
                dist[v] = nd
                heappush(pq, (nd, v))
    return dist
\end{lstlisting}

\subsection{Floyd}

\textbf{适用题型：}
点数较小的任意两点最短路，通常 $n \le 400$。

\begin{lstlisting}
def floyd(dist):
    n = len(dist)
    for k in range(n):
        for i in range(n):
            dik = dist[i][k]
            for j in range(n):
                if dik + dist[k][j] < dist[i][j]:
                    dist[i][j] = dik + dist[k][j]
    return dist
\end{lstlisting}

\subsection{最小生成树 Kruskal}

\begin{lstlisting}
def kruskal(n, edges):
    # edges = [(w, u, v), ...]
    dsu = DSU(n)
    edges.sort()
    total = 0
    cnt = 0
    for w, u, v in edges:
        if dsu.union(u, v):
            total += w
            cnt += 1
    return total if cnt == n - 1 else None
\end{lstlisting}

\section{动态规划}

\subsection{DP 设计四步}

\begin{enumerate}
  \item \textbf{定义状态：} $dp[i]$ 或 $dp[i][j]$ 表示什么。
  \item \textbf{确定转移：} 从哪些旧状态推到新状态。
  \item \textbf{初始化：} 边界状态是什么。
  \item \textbf{确定答案：} 是 $dp[n]$、最大值、最小值还是求和。
\end{enumerate}

\subsection{一维线性 DP}

\begin{lstlisting}
def max_subarray(a):
    best = cur = a[0]
    for x in a[1:]:
        cur = max(x, cur + x)
        best = max(best, cur)
    return best
\end{lstlisting}

\subsection{最长上升子序列}

\[
O(n^2) \quad \text{或} \quad O(n\log n)
\]

\begin{lstlisting}
from bisect import bisect_left

def lis_length(a):
    d = []
    for x in a:
        i = bisect_left(d, x)
        if i == len(d):
            d.append(x)
        else:
            d[i] = x
    return len(d)
\end{lstlisting}

\subsection{01 背包}

\textbf{适用题型：}
每个物品最多选一次。

\begin{lstlisting}
def zero_one_knapsack(items, W):
    # items = [(weight, value), ...]
    dp = [0] * (W + 1)
    for w, v in items:
        for cap in range(W, w - 1, -1):
            dp[cap] = max(dp[cap], dp[cap - w] + v)
    return dp[W]
\end{lstlisting}

\subsection{完全背包}

\textbf{适用题型：}
每个物品可以选无限次。

\begin{lstlisting}
def complete_knapsack(items, W):
    dp = [0] * (W + 1)
    for w, v in items:
        for cap in range(w, W + 1):
            dp[cap] = max(dp[cap], dp[cap - w] + v)
    return dp[W]
\end{lstlisting}

\subsection{区间 DP}

\textbf{适用题型：}
合并石子、括号匹配、区间划分。

\begin{lstlisting}
def interval_dp(a):
    n = len(a)
    INF = 10 ** 30
    dp = [[0] * n for _ in range(n)]
    pre = [0]
    for x in a:
        pre.append(pre[-1] + x)

    for length in range(2, n + 1):
        for l in range(n - length + 1):
            r = l + length - 1
            dp[l][r] = INF
            total = pre[r + 1] - pre[l]
            for k in range(l, r):
                dp[l][r] = min(dp[l][r], dp[l][k] + dp[k + 1][r] + total)
    return dp[0][n - 1]
\end{lstlisting}

\subsection{树形 DP}

\textbf{适用题型：}
树上最大独立集、树上路径、以子树为单位统计。

\begin{lstlisting}
import sys
sys.setrecursionlimit(10 ** 7)

def tree_dp(n, g, val):
    dp0 = [0] * n
    dp1 = [0] * n

    def dfs(u, p):
        dp1[u] = val[u]
        for v in g[u]:
            if v == p:
                continue
            dfs(v, u)
            dp0[u] += max(dp0[v], dp1[v])
            dp1[u] += dp0[v]

    dfs(0, -1)
    return max(dp0[0], dp1[0])
\end{lstlisting}

\subsection{状态压缩 DP}

\textbf{适用条件：}
$n$ 较小，一般 $n \le 20$。

\begin{lstlisting}
def tsp(dist):
    n = len(dist)
    INF = 10 ** 30
    dp = [[INF] * n for _ in range(1 << n)]
    dp[1][0] = 0
    for mask in range(1 << n):
        for u in range(n):
            if dp[mask][u] >= INF:
                continue
            for v in range(n):
                if mask >> v & 1:
                    continue
                nm = mask | (1 << v)
                dp[nm][v] = min(dp[nm][v], dp[mask][u] + dist[u][v])
    return min(dp[(1 << n) - 1])
\end{lstlisting}

\subsection{数位 DP}

\textbf{适用题型：}
统计 $[0,n]$ 内满足某种数字性质的数。

\begin{lstlisting}
from functools import lru_cache

def count_numbers(n):
    s = str(n)

    @lru_cache(None)
    def dfs(pos, tight, started, state):
        if pos == len(s):
            return 1 if started else 0

        limit = int(s[pos]) if tight else 9
        ans = 0
        for d in range(limit + 1):
            ntight = tight and (d == limit)
            nstarted = started or d != 0
            nstate = state
            if nstarted:
                # update nstate here
                pass
            ans += dfs(pos + 1, ntight, nstarted, nstate)
        return ans

    return dfs(0, True, False, 0)
\end{lstlisting}

\section{数学与数论}

\subsection{最大公约数与最小公倍数}

\begin{lstlisting}
from math import gcd

def lcm(a, b):
    return a // gcd(a, b) * b
\end{lstlisting}

\subsection{快速幂}

\begin{lstlisting}
def qpow(a, b, mod):
    res = 1
    a %= mod
    while b:
        if b & 1:
            res = res * a % mod
        a = a * a % mod
        b >>= 1
    return res
\end{lstlisting}

\subsection{逆元}

当模数 $MOD$ 是质数时：
\[
a^{-1}\equiv a^{MOD-2}\pmod{MOD}.
\]

\begin{lstlisting}
def inv(a, mod):
    return pow(a, mod - 2, mod)
\end{lstlisting}

\subsection{组合数预处理}

\begin{lstlisting}
def prepare_comb(n, mod):
    fac = [1] * (n + 1)
    ifac = [1] * (n + 1)
    for i in range(1, n + 1):
        fac[i] = fac[i - 1] * i % mod
    ifac[n] = pow(fac[n], mod - 2, mod)
    for i in range(n, 0, -1):
        ifac[i - 1] = ifac[i] * i % mod

    def C(n, k):
        if k < 0 or k > n:
            return 0
        return fac[n] * ifac[k] % mod * ifac[n - k] % mod

    return C
\end{lstlisting}

\subsection{素数筛}

\begin{lstlisting}
def sieve(n):
    is_prime = [True] * (n + 1)
    if n >= 0:
        is_prime[0] = False
    if n >= 1:
        is_prime[1] = False
    primes = []
    for i in range(2, n + 1):
        if is_prime[i]:
            primes.append(i)
            step = i
            start = i * i
            if start <= n:
                for j in range(start, n + 1, step):
                    is_prime[j] = False
    return primes, is_prime
\end{lstlisting}

\subsection{质因数分解}

\begin{lstlisting}
def factorize(x):
    res = []
    d = 2
    while d * d <= x:
        if x % d == 0:
            cnt = 0
            while x % d == 0:
                x //= d
                cnt += 1
            res.append((d, cnt))
        d += 1 if d == 2 else 2
    if x > 1:
        res.append((x, 1))
    return res
\end{lstlisting}

\subsection{矩阵快速幂}

\textbf{适用题型：}
线性递推、斐波那契变形、固定状态转移重复 $k$ 次。

\begin{lstlisting}
def mat_mul(A, B, mod):
    n, m, p = len(A), len(B), len(B[0])
    C = [[0] * p for _ in range(n)]
    for i in range(n):
        for k in range(m):
            if A[i][k]:
                aik = A[i][k]
                for j in range(p):
                    C[i][j] = (C[i][j] + aik * B[k][j]) % mod
    return C

def mat_pow(A, e, mod):
    n = len(A)
    R = [[0] * n for _ in range(n)]
    for i in range(n):
        R[i][i] = 1
    while e:
        if e & 1:
            R = mat_mul(R, A, mod)
        A = mat_mul(A, A, mod)
        e >>= 1
    return R
\end{lstlisting}

\section{字符串算法}

\subsection{KMP}

\textbf{适用题型：}
模式串匹配、求最短循环节、前后缀关系。

\begin{lstlisting}
def prefix_function(s):
    n = len(s)
    pi = [0] * n
    for i in range(1, n):
        j = pi[i - 1]
        while j and s[i] != s[j]:
            j = pi[j - 1]
        if s[i] == s[j]:
            j += 1
        pi[i] = j
    return pi

def kmp_find(text, pat):
    if not pat:
        return list(range(len(text) + 1))
    pi = prefix_function(pat)
    ans = []
    j = 0
    for i, ch in enumerate(text):
        while j and ch != pat[j]:
            j = pi[j - 1]
        if ch == pat[j]:
            j += 1
        if j == len(pat):
            ans.append(i - len(pat) + 1)
            j = pi[j - 1]
    return ans
\end{lstlisting}

\subsection{Trie 字典树}

\textbf{适用题型：}
前缀统计、字符串集合、异或最大值。

\begin{lstlisting}
class Trie:
    def __init__(self):
        self.ch = []
        self.end = []
        self.ch.append({})
        self.end.append(0)

    def insert(self, s):
        u = 0
        for c in s:
            if c not in self.ch[u]:
                self.ch[u][c] = len(self.ch)
                self.ch.append({})
                self.end.append(0)
            u = self.ch[u][c]
        self.end[u] += 1

    def search(self, s):
        u = 0
        for c in s:
            if c not in self.ch[u]:
                return 0
            u = self.ch[u][c]
        return self.end[u]
\end{lstlisting}

\subsection{字符串哈希}

\begin{lstlisting}
class StringHash:
    def __init__(self, s, base=911382323, mod=(1 << 61) - 1):
        self.mod = mod
        self.base = base
        n = len(s)
        self.h = [0] * (n + 1)
        self.p = [1] * (n + 1)
        for i, ch in enumerate(s):
            self.h[i + 1] = (self.h[i] * base + ord(ch)) % mod
            self.p[i + 1] = self.p[i] * base % mod

    def get(self, l, r):
        # [l, r)
        return (self.h[r] - self.h[l] * self.p[r - l]) % self.mod
\end{lstlisting}

\subsection{Manacher}

\textbf{适用题型：}
最长回文子串、统计回文半径。

\begin{lstlisting}
def manacher(s):
    t = "#" + "#".join(s) + "#"
    n = len(t)
    p = [0] * n
    center = right = 0
    best = 0
    for i in range(n):
        mirror = 2 * center - i
        if i < right:
            p[i] = min(right - i, p[mirror])
        while i - p[i] - 1 >= 0 and i + p[i] + 1 < n and t[i - p[i] - 1] == t[i + p[i] + 1]:
            p[i] += 1
        if i + p[i] > right:
            center, right = i, i + p[i]
        best = max(best, p[i])
    return best
\end{lstlisting}

\section{搜索、回溯与剪枝}

\subsection{回溯模板}

\textbf{适用题型：}
排列、组合、棋盘、子集枚举、小规模构造。

\begin{lstlisting}
def backtrack(path, used, nums, ans):
    if len(path) == len(nums):
        ans.append(path[:])
        return
    for i, x in enumerate(nums):
        if used[i]:
            continue
        used[i] = True
        path.append(x)
        backtrack(path, used, nums, ans)
        path.pop()
        used[i] = False
\end{lstlisting}

\subsection{Meet in the Middle}

\textbf{适用条件：}
$n$ 大约 $30$ 到 $44$，直接 $2^n$ 不可行，但 $2^{n/2}$ 可行。

\begin{lstlisting}
from bisect import bisect_right

def subset_sums(a):
    res = [0]
    for x in a:
        res += [y + x for y in res]
    return res

def count_subset_sum_leq(a, limit):
    mid = len(a) // 2
    left = subset_sums(a[:mid])
    right = subset_sums(a[mid:])
    right.sort()
    ans = 0
    for x in left:
        ans += bisect_right(right, limit - x)
    return ans
\end{lstlisting}

\section{位运算}

\subsection{常用技巧}

\begin{longtable}{p{0.32\textwidth}p{0.55\textwidth}}
\toprule
\textbf{表达式} & \textbf{含义} \\
\midrule
\texttt{x \& 1} & 判断奇偶 \\
\texttt{x >> k \& 1} & 判断第 $k$ 位是否为 1 \\
\texttt{x | (1 << k)} & 把第 $k$ 位置 1 \\
\texttt{x \& \textasciitilde(1 << k)} & 把第 $k$ 位置 0 \\
\texttt{x \^{} (1 << k)} & 翻转第 $k$ 位 \\
\texttt{x \& -x} & 取最低位的 1 \\
\texttt{x.bit\_count()} & 统计 1 的个数 \\
\bottomrule
\end{longtable}

\subsection{枚举子集}

\begin{lstlisting}
def enumerate_submasks(mask):
    sub = mask
    res = []
    while sub:
        res.append(sub)
        sub = (sub - 1) & mask
    res.append(0)
    return res
\end{lstlisting}

\section{扫描线与区间覆盖}

\subsection{最大重叠区间数量}

\textbf{适用题型：}
会议室数量、同时在线人数、区间覆盖层数。

\begin{lstlisting}
def max_overlap(intervals):
    events = []
    for l, r in intervals:
        events.append((l, 1))
        events.append((r, -1))
    events.sort(key=lambda x: (x[0], x[1]))
    cur = ans = 0
    for _, delta in events:
        cur += delta
        ans = max(ans, cur)
    return ans
\end{lstlisting}

\subsection{整数闭区间覆盖}

若区间是闭区间 $[l,r]$，可把结束事件放在 $r+1$。

\begin{lstlisting}
def max_overlap_closed(intervals):
    diff = []
    for l, r in intervals:
        diff.append((l, 1))
        diff.append((r + 1, -1))
    diff.sort()
    cur = ans = 0
    for _, d in diff:
        cur += d
        ans = max(ans, cur)
    return ans
\end{lstlisting}

\section{博弈论}

\subsection{Nim 游戏}

若若干堆石子，每次选择一堆拿任意正数个，则：
\[
a_1 \oplus a_2 \oplus \cdots \oplus a_n \ne 0
\]
为先手必胜。

\begin{lstlisting}
def nim_winner(piles):
    x = 0
    for v in piles:
        x ^= v
    return "First" if x else "Second"
\end{lstlisting}

\subsection{SG 函数}

\begin{lstlisting}
def mex(s):
    g = 0
    while g in s:
        g += 1
    return g

def sg_values(n, moves):
    sg = [0] * (n + 1)
    for i in range(1, n + 1):
        reach = set()
        for m in moves:
            if i >= m:
                reach.add(sg[i - m])
        sg[i] = mex(reach)
    return sg
\end{lstlisting}

\section{常见综合题型}

\subsection{最值 + 可行性}

\textbf{题型：}
给一个答案 $x$，问能否在限制内完成。

\textbf{套路：}
二分答案 + 贪心/前缀/图搜索判断。

\[
\text{答案单调} \Rightarrow \text{二分}
\]

\subsection{区间修改 + 区间查询}

\textbf{题型：}
多次修改数组区间，并查询区间和/最值。

\textbf{套路：}
\begin{itemize}
  \item 只有最终结果：差分。
  \item 单点修改 + 区间和：树状数组。
  \item 区间修改 + 区间查询：线段树。
\end{itemize}

\subsection{连通块问题}

\textbf{题型：}
网格岛屿、朋友关系、加边后连通性。

\textbf{套路：}
\begin{itemize}
  \item 静态图：DFS/BFS。
  \item 动态合并：并查集。
  \item 带权边最小连接：Kruskal。
\end{itemize}

\subsection{计数取模}

\textbf{题型：}
方案数、排列数、路径数。

\textbf{套路：}
\begin{itemize}
  \item 有明显递推：DP。
  \item 选择若干元素：组合数。
  \item 大指数：快速幂。
  \item 线性递推次数很大：矩阵快速幂。
\end{itemize}

\subsection{字符串匹配}

\textbf{题型：}
判断模式串出现次数、最长公共前缀、回文。

\textbf{套路：}
\begin{itemize}
  \item 单模式匹配：KMP。
  \item 多字符串前缀：Trie。
  \item 子串快速比较：字符串哈希。
  \item 回文子串：Manacher。
\end{itemize}

\section{比赛实用技巧}

\subsection{读题顺序}

\begin{enumerate}
  \item 先看数据范围。
  \item 再看是否多组测试。
  \item 找关键词：区间、路径、取模、最小最大、方案数。
  \item 先写暴力，判断瓶颈。
  \item 用对应算法优化瓶颈。
\end{enumerate}

\subsection{Python 常见坑}

\begin{itemize}
  \item 递归深度：树 DFS 前设置 \texttt{sys.setrecursionlimit}。
  \item 大输入：使用 \texttt{sys.stdin.readline}。
  \item 取模：每次乘法或加法后及时取模。
  \item 浮点二分：控制迭代次数，如 80 次。
  \item 多组测试：注意清空全局数组。
  \item 下标：明确 0-index 还是 1-index。
\end{itemize}

\subsection{浮点二分模板}

\begin{lstlisting}
def float_binary(lo, hi, check):
    for _ in range(80):
        mid = (lo + hi) / 2
        if check(mid):
            hi = mid
        else:
            lo = mid
    return (lo + hi) / 2
\end{lstlisting}

\subsection{常用取模写法}

\begin{lstlisting}
MOD = 10 ** 9 + 7

x = (x + y) % MOD
x = x * y % MOD
x = (x - y) % MOD
\end{lstlisting}

\section{赛前复习清单}

\begin{itemize}
  \item 会写快速输入模板。
  \item 会写二分答案，并能判断单调性。
  \item 会写前缀和、差分、离散化。
  \item 会写并查集、树状数组、线段树。
  \item 会写 BFS、Dijkstra、拓扑排序。
  \item 会写 01 背包、完全背包、LIS、区间 DP。
  \item 会写快速幂、组合数、素数筛。
  \item 会写 KMP、Trie、字符串哈希。
  \item 熟悉常见边界：空数组、单元素、负数、闭区间、取模。
\end{itemize}

\section{最后的选型口诀}

\begin{itemize}
  \item 看到“多次区间求和”，先想前缀和。
  \item 看到“多次区间修改”，先想差分或线段树。
  \item 看到“最大最小/最小最大”，先想二分答案。
  \item 看到“连通性”，先想并查集或 DFS。
  \item 看到“最短路”，先看边权：无权 BFS，非负 Dijkstra。
  \item 看到“方案数取模”，先想 DP 或组合数学。
  \item 看到“字符串匹配”，先想 KMP 或哈希。
  \item 看到“坐标巨大但数量少”，先想离散化。
  \item 看到“$n \le 20$”，先想状态压缩。
  \item 看到“两人轮流操作”，先想博弈论。
\end{itemize}

\end{document}